%!TEX root = ../../thesis.tex

\section{Reinforcement Learning}\label{sec:reinforcement-learning}

% Drafting outline: bullet points to be fleshed out into prose.

\subsection{Agent, Environment, and Markov Decision Processes}\label{sec:rl-mdp}

% --- Verbatim excerpt from the proposal (own text; citation keys remapped) ---
Reinforcement learning (RL) is an extremely general framework for approaching problems
in which an agent strives to maximize a reward signal through a selection of actions
in certain states~\cite{sutton_reinforcement_2020}.
These kinds of problems can be modeled as finite Markov decision processes (MDPs).
A Markov decision process is a 3-tuple $MDP=(S,A,p)$ where

\begin{itemize}
  \item $\mathcal{S}$ is the set of states,
  \item $\mathcal{A}$ is the set of actions where $\mathcal{A}_s\subseteq A$ are actions available in state $s\in\mathcal{S}$,
  \item $p(s^\prime,r|s,a)=\Pr\{S_t=s^\prime, R_t=r|S_{t-1}=s, A_{t-1}=a\}$
  describes the system dynamics, i.e., the probability of encountering a next
  state $s^\prime$ and reward $r$ at time step $t$ when taking action $a$ in
  state $s$ at time step $t-1$
\end{itemize}

An RL agent's goal is to learn a policy $\pi:\mathcal{S}\times\mathcal{A}\rightarrow [0,1]$
mapping each state $s$ to probabilities of taking each possible action $a$ to maximize
the cumulative reward.
% --- End excerpt (notation cleanup when fleshing out: unify $S$ vs.\ $\mathcal{S}$
%     in the tuple; consider extending to $(\mathcal{S},\mathcal{A},p,\gamma)$) ---

\begin{itemize}
  \item Distinguish from supervised learning: no labeled correct actions, only evaluative
    (and possibly delayed) feedback; trial-and-error credit assignment
    ~\cite{sutton_reinforcement_2020}; idea traces back to early AI~\cite{minsky_steps_1961};
    add \cite{puterman_markov_2009} as the standard MDP reference next to the excerpt.
  \item Extend the excerpt's formalism~\cite{sutton_reinforcement_2020}:
    \begin{itemize}
      \item Markov property: transition depends only on current state and action,
      \item episodic vs.\ continuing tasks; discount factor $\gamma$ and discounted return
        $G_t = \sum_k \gamma^k r_{t+k}$,
      \item state value $V^\pi$, action value $Q^\pi$; Bellman (optimality) equations;
        optimal policy $\pi^*$.
    \end{itemize}
  \item Partial observability: agent often sees an \emph{observation}, not the full state
    (POMDP setting)~\cite{sutton_reinforcement_2020} --- relevant later: a coverage vector is a
    very partial view of the emulated program's true state
    (\cref{chap:efficientzero-fuzzing}).
\end{itemize}

\subsection{Exploration versus Exploitation}\label{sec:rl-exploration}

\begin{itemize}
  \item Fundamental dilemma: exploit the best-known action or explore to find better
    ones~\cite{sutton_reinforcement_2020}.
  \item Minimal setting: multi-armed bandit (single state); $\varepsilon$-greedy and
    optimism-in-the-face-of-uncertainty / UCB-style strategies
    ~\cite{sutton_reinforcement_2020,hutchison_bandit_2006}.
  \item Contextual bandit as intermediate step between bandit and full MDP ---
    mirrors the ladder of evaluation targets in \cref{chap:targets}
    (\cref{sec:target-bandit,sec:target-contextual-bandit,sec:target-sequence}).
  \item Foreshadow: fuzzing has exactly this structure --- keep hammering a productive protocol
    state vs.\ probe untested ones~\cite{borcherding_bandits_2023}.
\end{itemize}

\subsection{Algorithm Families}\label{sec:rl-algorithms}

% --- Verbatim excerpt from the proposal (own text; citation keys remapped).
%     Key changes vs.\ proposal: williams1992simple (REINFORCE) is not in references/*.bib yet
%     (todo below); mnih2013playing (arXiv DQN) mapped to the Nature version
%     mnih_human-level_2015; hafner2023mastering mapped to hafner_mastering_2024.
%     Small fixes for fleshing out: "make use modern-day" -> "make use of modern-day";
%     $Q:\mathcal{S}\rightarrow\mathcal{R}$ should arguably be
%     $Q:\mathcal{S}\times\mathcal{A}\rightarrow\mathbb{R}$. ---
Algorithms for finding an optimal $\pi^*$ can be divided into three classes:
value-based, policy-based and actor-critic.

Value-based methods try to learn a value function $Q:\mathcal{S}\rightarrow \mathcal{R}$
mapping states to their expected cumulative reward so that in each step, the
action that maximizes the expected reward can be taken.
These include well-known methods such as Q-learning, TD($\lambda$) or SARSA.

Policy-based methods try to directly construct an optimal policy $\pi^*$ without the
use of a value function.
They have various advantages.
For example, they can easily handle continuous action spaces and ensure stochasticity
which is important for state space exploration.
A well-known classical policy-based algorithm is
REINFORCE\todo{cite: Williams 1992, ``Simple statistical gradient-following algorithms
for connectionist reinforcement learning'' (was williams1992simple in the proposal)}.

Actor-critic methods synthesize value- and policy-based methods by including information
about a state's value in the policy learning process.

Since classical, tabular methods don't work well with large state spaces,
state-of-the-art methods make use modern-day deep learning advancements.
Examples include the family of Deep-Q Network (DQN)
algorithms~\cite{mnih_human-level_2015,hessel_rainbow_2017} approximating the $Q$
function by neural networks, policy-based gradient-descent methods like
Proximal Policy Optimization (PPO)~\cite{schulman_proximal_2017} and actor-critic
methods like A3C~\cite{mnih_asynchronous_2016}.
Impressive results range from the master-level Go program
AlphaGo~\cite{silver_mastering_2016} and its more generalized successors
AlphaZero~\cite{silver_mastering_2017} and MuZero~\cite{schrittwieser_mastering_2020}
to Dreamerv3~\cite{hafner_mastering_2024} that is able to learn how to acquire
diamonds in the popular video game \textit{Minecraft} without any human input
or expert knowledge.
% --- End excerpt ---

% Complementary bullets not covered by the excerpt:
\begin{itemize}
  \item Foundational citations to attach to the excerpt's method names:
    temporal-difference learning~\cite{sutton_learning_1988}; Q-learning
    \todo{cite: Watkins \& Dayan 1992, ``Q-learning''}.
  \item Further DQN refinements between DQN and Rainbow: double DQN~\cite{hasselt_deep_2015},
    dueling networks~\cite{wang_dueling_2016}, prioritized experience
    replay~\cite{schaul_prioritized_2016} (PER returns in EfficientZero,
    \cref{chap:efficientzero}).
  \item Further policy-gradient/actor-critic methods: TRPO~\cite{schulman_trust_2017},
    SAC~\cite{haarnoja_soft_2018}.
  \item Broad algorithm surveys~\cite{ghasemi_comprehensive_2025}.
  \item Practical caveats of deep RL: high variance across seeds, hyperparameter sensitivity,
    reproducibility concerns~\cite{henderson_deep_2019} --- informs the evaluation methodology in
    \cref{sec:analysis-evaluation}.
\end{itemize}

\subsection{Model-Free versus Model-Based RL}\label{sec:rl-model-based}

\begin{itemize}
  \item Model-free: learn policy/values directly from transitions; conceptually simple, but
    notoriously sample-hungry --- millions to billions of environment steps
    ~\cite{mnih_human-level_2015,moerland_model-based_2022}.
  \item Model-based: additionally learn a \emph{dynamics model} of the environment and use it for
    planning or imagined training --- the main lever for sample
    efficiency~\cite{moerland_model-based_2022,kaiser_model-based_2024}.
  \item World-model line: learn compressed latent dynamics, train the policy inside the
    model~\cite{ha_world_2018,hafner_learning_2019,hafner_mastering_2024}.
  \item Planning-with-search line:
    \begin{itemize}
      \item Monte-Carlo tree search (MCTS)~\cite{hutchison_efficient_2007}, UCT: bandit-based
        selection inside the tree~\cite{hutchison_bandit_2006}, modern
        review~\cite{swiechowski_monte_2023},
      \item AlphaGo: MCTS + learned value/policy networks, with known rules
        ~\cite{silver_mastering_2016}; AlphaGo Zero: self-play only
        ~\cite{silver_mastering_2017-1}; AlphaZero: generalized across
        games~\cite{silver_general_2018},
      \item MuZero: replaces the known simulator with a \emph{learned latent model} --- planning
        without knowing the rules~\cite{schrittwieser_mastering_2020},
      \item EfficientZero: MuZero made sample-efficient (self-supervised consistency, value
        prefix, model-based off-policy correction); strong Atari performance from only 100k
        steps ($\approx$ 2 hours) of experience~\cite{ye_mastering_2021} --- treated in depth in
        \cref{chap:efficientzero}.
    \end{itemize}
\end{itemize}

\subsection{Sample Efficiency}\label{sec:rl-sample-efficiency}

\begin{itemize}
  \item Define: performance achieved per environment interaction; the binding constraint whenever
    a step is expensive~\cite{moerland_model-based_2022,ye_mastering_2021}.
  \item In this thesis each environment step is a program execution under emulation
    (\cref{sec:fuzzing-emulation}) --- orders of magnitude slower than a simulator frame, so the
    fuzzing use case sits squarely in the low-sample regime that EfficientZero targets
    ~\cite{ye_mastering_2021}.
  \item This is the central argument for choosing a model-based, planning-heavy algorithm over the
    model-free methods used by most prior RL fuzzers (\cref{sec:reinforcement-fuzzing}).
\end{itemize}
